% =====================================================================
%  draft.tex — Thesis DRAFT skeleton for the Score-Bundle Models project.
%
%  This is a scaffold, not a finished thesis: prose is written where the
%  result is known and honest, and \TODO{...} marks everything to fill in
%  (numbers, figures, citations, expansions).  Exact figures live in
%  docs/phase1_calibration_results.md, docs/downstream_tasks_results.md,
%  docs/status_2026-07-02.md; the full model derivation is docs/formulation.tex.
%
%  Compiles with plain pdflatex (no bibtex needed — manual thebibliography).
%  Numbers are centralised in the macros below so there is ONE place to edit.
% =====================================================================
\documentclass[11pt,a4paper]{report}

\usepackage[utf8]{inputenc}
\usepackage[T1]{fontenc}
\usepackage[margin=1in]{geometry}
\usepackage{amsmath,amssymb,amsthm}
\usepackage{graphicx}
\usepackage{booktabs}
\usepackage{xcolor}
\usepackage{enumitem}
\usepackage[hidelinks]{hyperref}

\graphicspath{{figures/}}

% --- draft helpers ---------------------------------------------------
\newcommand{\TODO}[1]{{\color{red}\textbf{[TODO: #1]}}}
\newcommand{\placeholder}[1]{{\color{gray}\textit{#1}}}
% Figure placeholder: shows the figure if present, else a TODO box.
\newcommand{\figph}[3]{% {file}{width}{caption}
  \begin{figure}[t]\centering
  \IfFileExists{figures/#1}{\includegraphics[width=#2\linewidth]{#1}}%
    {\fbox{\parbox[c][0.28\linewidth][c]{#2\linewidth}{\centering\TODO{figure: #1}}}}
  \caption{#3}\end{figure}}

% --- headline numbers (UPDATE from docs/phase1_calibration_results.md) ---
\newcommand{\rmseZero}{0.566}
\newcommand{\rmseGraph}{0.404}
\newcommand{\rmseNetGraph}{0.393}      % published headline (strict protocol)
\newcommand{\rmseStack}{0.388}         % features + network + graph (candidate, strict)
\newcommand{\nllZero}{$-0.007$}
\newcommand{\nllGraph}{$-0.308$}
\newcommand{\nllNetGraph}{$-0.322$}    % strict protocol
\newcommand{\nllStack}{$-0.333$}       % strict protocol
\newcommand{\covZero}{87\%}
\newcommand{\covHead}{92\%}            % coverage of nominal-90\% intervals
\newcommand{\nPieces}{30}
\newcommand{\maskRate}{40\%}

% --- notation (must match docs/formulation.tex and CLAUDE.md) ---------
\newcommand{\Sy}{\Sigma_y}             % posterior covariance
\newcommand{\muLM}{\mu_{\mathrm{LM}}}  % LM prior mean

\title{\textbf{Score-Bundle Models}\\[0.3em]
  \large Bayesian Score-Informed Performance Transcription\\
  with Calibrated Uncertainty \TODO{finalise title}}
\author{Ray \TODO{full name} \\[0.5em]
  \small Supervisor: \TODO{name} \\
  \small \TODO{institution / degree programme}}
\date{Draft --- \today}

\begin{document}
\maketitle

\begin{abstract}
Given a printed score and a recording of a performance (as piano MIDI in this
work), we recover \emph{how} each written note was played --- its micro-timing,
articulation, and dynamics --- together with a calibrated statement of
confidence. The core idea is a \emph{score graph}: notes that are related in the
score are modelled as related in performance, via a Gaussian Markov random field
(GMRF) prior over per-note expressive variables. A small, from-scratch symbolic
music language model supplies a learned prior mean, and the graph models the
residual with structured, calibrated uncertainty. On held-out ASAP piano
performances the full model recovers hidden notes with lower error than any
per-note predictor (RMSE~$\approx\rmseNetGraph$) and better-calibrated confidence
(nominal-90\% intervals cover the truth $\approx\covHead$ of the time), with the
graph's contribution statistically significant on both axes. Six downstream tasks
map exactly where the structure helps (per-note recovery, confidence, error
spotting) and where it does not (whole-performance summaries).
\TODO{tighten abstract; insert final numbers and the one headline sentence.}
\end{abstract}

\tableofcontents

% =====================================================================
\chapter{Introduction}\label{ch:intro}

\section{Motivation}
\TODO{Why expressive performance matters; the gap between a score (what to play)
and a performance (how it was played); why uncertainty is essential.}
A score specifies \emph{what} to play; a performance realises \emph{how}. The
difference --- micro-timing, articulation, dynamics --- is where musical
expression lives, and recovering it from a recording is both scientifically
interesting and practically useful (analysis, pedagogy, expressive rendering).

\section{Problem statement}
We study \emph{score-informed} performance transcription: the symbolic score
(note identities, onsets, durations) is \emph{known}, and the task is to infer
the continuous performance variables attached to each note, with uncertainty.
This is the inverse of expressive synthesis. Inferring the score itself
(full automatic music transcription) is out of scope; see \S\ref{sec:scope}.
\TODO{formal one-paragraph statement; forward/inverse framing.}

\section{Contributions}
\begin{enumerate}[leftmargin=*]
  \item A \textbf{score-graph structured prior} (GMRF/SPDE) over per-note
        expressive variables, giving a closed-form posterior with calibrated
        per-note uncertainty, validated on held-out data.
  \item An \textbf{honest characterisation of a learned prior mean}: a
        from-scratch symbolic music LM contributes calibration- and
        loudness-specific information; against a rich hand-built score-feature
        baseline it \emph{ties on average error} and the two \emph{stack}.
  \item A \textbf{map of where the structure helps}, via six downstream tasks
        with objective metrics and honest baselines (three wins, two mixed, one
        negative).
  \item \textbf{Methodological rigour}: a discovered-and-fixed evaluation leak,
        a strict leakage-free protocol, and a robustness guard for the
        hyperparameter fit. \TODO{phrase as a contribution, not an apology.}
  \item \placeholder{(Planned / future)} A path to audio via a differentiable
        synthesizer likelihood (Phase~3) and a geometric ``connection-Laplacian''
        generalisation of the graph prior.
\end{enumerate}

\section{Scope and non-goals}\label{sec:scope}
In scope: piano, symbolic (MIDI) performance variables (timing, articulation,
dynamics), known score. Not (yet) in scope: audio input/output, intonation and
vibrato (require continuously pitched instruments), and unknown-score inference.
\TODO{expand; state assumptions explicitly.}

\section{Thesis outline}
\TODO{one line per chapter.}

% =====================================================================
\chapter{Background and Related Work}\label{ch:related}

\section{Computational models of expressive performance}
\TODO{survey; position our per-note, uncertainty-first framing.}
\cite{cancino2018,oore2018}

\section{Score-to-performance alignment}
\TODO{how aligned corpora are built; ASAP; note-level alignment.} \cite{asap}

\section{Structured priors: GMRFs and the SPDE approach}
\TODO{graph Laplacian priors; the SPDE/GMRF link; Mat\'ern GPs on graphs.}
\cite{lindgren2011,rueheld2005,borovitskiy2021}

\section{Symbolic music language models}
\TODO{Music Transformer, REMI, performance RNNs; from-scratch (nanoGPT) rationale.}
\cite{musictransformer,remi,oore2018,nanogpt}

\section{Differentiable synthesis}
\TODO{DDSP; harmonic-plus-noise; why it matters for Phase~3.} \cite{ddsp}

% =====================================================================
\chapter{Model and Formulation}\label{ch:model}
\emph{The complete derivation is in \texttt{docs/formulation.tex}
(Appendix~\ref{app:formulation}); the essentials are reproduced here.}

\section{Score and performance variables}
The known score is $\mathcal{S}=\{s_i\}_{i=1}^N$ with $s_i=(p_i,b_i,d_i)$ (pitch as
MIDI number, beat onset, beat duration). Each note carries a Phase-1 performance
vector
\[
  y_i = [\,\tau_i,\ \log r_i,\ v_i\,]\in\mathbb{R}^3,
\]
the onset residual (s), the articulation/duration ratio (log), and the
(normalised) velocity. \TODO{define $\tau$ relative to the estimated tempo warp.}

\section{The score graph}
Build a weighted graph $G=(V,E)$ over notes with edge weights from score-time,
pitch, voice, and phrase adjacency, with Laplacian $L_G=D-W$. \TODO{final edge
definition; which adjacencies were used.}

\section{Graph prior (GMRF)}
Place a Gaussian prior with precision built from the graph,
\[
  Q_G = \lambda I + \eta L_G
  \qquad\text{(or Mat\'ern/SPDE }Q_G=\sigma_g^{-2}(\kappa^2 I + L_G)^{\alpha}),
\]
so that connected notes have correlated expressive deviations, with a closed-form
Gaussian posterior. \TODO{state per-channel construction.}

\section{Observation and posterior}
With $\tilde y = y + e,\ e\sim\mathcal N(0,\Sigma_e)$ on the \emph{observed} notes, write
$P$ for the diagonal observation precision --- $\Sigma_e^{-1}$ at observed nodes, $0$ at
held-out nodes, so held-out notes carry no likelihood term. The posterior is Gaussian,
\[
  p(y\mid\tilde y)=\mathcal N(m,\ \Sy),\quad
  \Sy=(Q_G+P)^{-1},\quad
  m=\Sy(Q_G\,\mu+P\,\tilde y),
\]
and held-out notes are predicted from their neighbours through the conditional Gaussian.
A held-out \emph{observation} is $y_i + e_i$, so its predictive variance must include the
observation noise, $(\Sy)_{ii} + \sigma_e^2$; scoring against the latent variance
$(\Sy)_{ii}$ alone is catastrophically overconfident (NLL and coverage blow up).

\section{Learned prior mean and the leak-free read-out}
The prior mean $\mu$ can be the LM prediction $\muLM$: a from-scratch decoder-only
Transformer, pretrained by next-token prediction on MAESTRO, read out per note.
\textbf{Crucially}, the read-out is taken \emph{before} a note's own performed
velocity enters the network, so it cannot leak the target.
\TODO{Fig.: the 4-token group and the pre-velocity read-out position.}

\section{Hyperparameters, fit, and robustness}
Prior hyperparameters are learned by maximizing the marginal likelihood of the
observed notes, per piece and per channel.
Two safeguards: a floor on the observation-noise level inside the fit, and a
self-validating guard that falls back to a conservative fit when the fitted field
fails on a held-back sliver of observed notes (precise statements in
Appendix~\ref{app:formulation}, \S5).

% =====================================================================
\chapter{Data and Experimental Protocol}\label{ch:data}

\section{Pretraining corpus (MAESTRO)}
Pretraining uses only the MAESTRO corpus \cite{maestro} (piano performances
captured on Disklavier instruments, as MIDI). \TODO{size, split, tokenisation.}

\section{Evaluation corpus (ASAP) and target extraction}
ASAP is the only corpus pairing written scores with performances; it supplies the
targets $y=[\tau,\log r,v]$ via the score$\leftrightarrow$performance beat warp.
\TODO{alignment/matching procedure; how many pieces/notes.}

\section{Contamination control and the strict protocol}
Because ASAP reuses many MAESTRO recordings, every ASAP performance whose MAESTRO
twin was in pretraining is removed (1036 $\to$ 653 aligned performances), so the
model is tested only on unseen performances. A \emph{strict} protocol
additionally blanks every hidden note's loudness from the input; its measured cost
is essentially nil on error (+0.0003 for the network system, +0.0000 for the
feature--network stack) and a small, honest calibration price ($\approx$+0.014
NLL) --- the strict numbers are the ones quoted throughout.

\section{Metrics}
Recovery: RMSE (per channel and pooled). Calibration: negative log predictive
density (NLL), central-interval coverage, PIT. \TODO{define; report median
alongside mean (single-cell collapse guard).}

\section{Performer-identification corpus (Vienna 4x22)}
For the one task ASAP cannot support (performer identity), the small Vienna~4x22
corpus (22 pianists $\times$ 4 shared excerpts) is used. \cite{vienna4x22}

% =====================================================================
\chapter{Phase-1 Results}\label{ch:results}

\section{Headline: recovery and calibration}
On \nPieces{} held-out ASAP pieces with \maskRate{} of notes hidden, the full
model (LM guess + score graph) attains the lowest error and best-calibrated
confidence; the graph's contribution is significant on \emph{both} axes.

\begin{table}[t]\centering
\caption{Held-out ASAP, strict protocol (numbers verified against
docs/phase1\_calibration\_results.md, 2026-07-03). Lower RMSE/NLL is better; coverage
closer to 90\% is better. \TODO{add per-piece bootstrap CIs.}}
\begin{tabular}{lccc}
\toprule
System & RMSE & NLL & coverage \\
\midrule
Predict zero                 & \rmseZero     & \nllZero  & \covZero \\
Graph alone                  & \rmseGraph    & \nllGraph & \covHead \\
Network + graph (headline)   & \textbf{\rmseNetGraph} & {\boldmath\nllNetGraph} & \covHead \\
Features + network + graph   & \rmseStack    & \nllStack & \covHead \\
\bottomrule
\end{tabular}
\end{table}

\figph{digest_headline.png}{0.7}{The graph's contribution per prior mean.
\TODO{final figure.}}

\section{The evaluation leak: discovery, fix, re-measurement}
An earlier read-out let the network see a note's own loudness --- the very target
it was asked to predict --- inflating the apparent loudness advantage. We read the
summary one position earlier (provably blind to the note's own loudness, pinned by
an automated test), re-measured everything, and adopted the strict protocol. The
corrected advantage is smaller but honest and statistically significant: on the
loudness channel, the network's mean-only error was 0.057 \emph{with} the leak (the
retracted claim), 0.118 under a first band-aid fix that also destroyed neighbour
context, and 0.090 with the final leak-free read-out.

\section{Hyperparameter fragility and the guard}
Two distinct failure modes were found and fixed. First, on a minority of mask
realisations the fitted noise level collapsed toward zero (overconfident
intervals); a noise-level floor inside the fit removes it. Second, a rarer
knife-edge collapse of the fitted coupling itself: one evaluation cell in 120 ---
always the same piece --- produced timing errors of 4.5 and 17 in two independent
runs, single-handedly dominating the pooled average (e.g.\ pooled timing RMSE
0.16 $\to$ 1.58). The self-validating guard removes both catastrophic cells
(4.5~$\to$~0.24; 17~$\to$~healthy) while being bit-for-bit identical to the
unguarded fit on the published protocol --- it never fires on a healthy fit.
Pooled tables additionally report the median and worst per-cell error, so any
future collapse is visible by construction.
\figph{digest_collapse.png}{0.7}{The 1-in-120 fit collapse, guard off vs.\ on.}

\section{Do hand-built features suffice?}
Against a rich 25-feature score baseline fit under the same protocol, the network
\emph{ties on average error}; its genuine, significant edge is calibration and
loudness. Best of all, features and network \emph{stack}
(RMSE $\rmseStack$ vs.\ $\rmseNetGraph$). \figph{digest_channels.png}{0.7}{Where
the network earns its place (per channel). \TODO{final figure.}}

\section{Stage-2: task-matched pretraining, an honest negative}
A bidirectional, mask-aware retraining (matched exactly to the evaluation task)
did \emph{not} beat the original at matched compute: pooled error is slightly but
significantly worse (+0.011), confidence quality ties strict-to-strict
($-0.325$ vs.\ $-0.322$), and a 3$\times$ budget does not close the gap (timing
improves, articulation and loudness do not). The honest lesson: the original
read-out already extracted what the objective change was meant to unlock. En
route, the same leak pattern tried to re-enter through the new bidirectional
read-out (a note read at its own velocity token \emph{contains} that velocity);
it was caught by its signature and fixed with a leave-one-out read-out, now a
documented design rule pinned by a structural test. \TODO{full table.}

\section{Scaling: a bigger LM does not help}
A 1.3$\times$ larger, lower-perplexity model (25.6M parameters, validation
perplexity 9.66 vs.\ 10.85) gives an essentially identical downstream result once
the leak is removed (RMSE 0.394 vs.\ 0.394, identical within noise) --- the
earlier ``scaling'' gain was the loudness leak. Practical rule: judge the LM
downstream, not by perplexity. (Measured with the earlier score-only read-out;
conservative for the LM rows.)

% =====================================================================
\chapter{Downstream Tasks}\label{ch:downstream}
Six tasks with objective metrics and honest baselines; all outcomes reported.
\TODO{one results table per task; move detail from docs/downstream\_tasks\_results.md.}

\section{Error spotting (anomaly detection) --- clear win}
Corrupt a small fraction of notes; rank by suspiciousness. The graph model ranks
better on every channel (AUROC 0.94~$\to$~0.98 at 3$\sigma$ corruptions,
0.85~$\to$~0.94 at 2$\sigma$) and its edge \emph{grows} as errors get subtler.

\section{Denoising --- win at known noise level}
Recover true values from added noise. With the noise level known, the graph beats
per-note shrinkage on both error and calibration (all 8 contrasts significant, at
matched $\approx$0.90 coverage). \emph{Honest negative}: estimating the noise
level blind from a single piece is unreliable --- intervals collapse to 0.54--0.64
coverage --- so transcription noise must be treated as known per system.

\section{Completion / rendering from an excerpt --- partial}
Predict the expression of the rest of a piece from its opening. Interpolation
(scattered hidden notes) is the clear win --- RMSE 0.450~$\to$~0.393, NLL
0.21~$\to$~0.01 at half the notes observed. Extrapolating far past the excerpt,
the LM guess carries the prediction and the graph is neutral or slightly harmful:
the graph helps only where hidden notes have nearby heard neighbours.

\section{Selective prediction --- qualified win}
Discarding the least-confident predictions lowers error on the rest, most clearly
for articulation (keeping the confident half cuts its RMSE 0.83~$\to$~0.54) and
for catching the model's own misfires (a timing blow-up at 0.61 overall falls to
0.07 on the confident half). Only the graph provides the per-note confidence
this triage needs; the no-graph baselines rank at piece level only.

\section{Musical-era classification --- honest negative}
Style statistics alone do not beat the majority-class baseline (accuracy
0.37--0.38 vs.\ 0.489 majority); the graph, a note-level tool, does not help
whole-piece summaries.

\section{Performer identification (Vienna 4x22) --- mixed}
The extracted expression \emph{does} carry a pianist's identity --- accuracy
0.136--0.179 against a 1-of-22 chance rate of 0.045 (3--4$\times$ chance),
validating the extraction. Graph smoothing \emph{hurts} this whole-performance
summary (0.080): denoising toward the ensemble mean removes exactly the
idiosyncrasy that identifies a performer.

\section{The boundary}
Across all six: \textbf{the graph prior helps whenever notes are judged one by one
for accuracy or confidence, and does not help when a performance is collapsed into
a summary.} This is exactly the thesis claim, now bounded by evidence on both
sides.

% =====================================================================
\chapter{Discussion}\label{ch:discussion}
\section{Interpretation}
\TODO{why structure helps calibration; why the LM's edge is loudness.}
\section{Limitations}
\TODO{symbolic-only; piano-only; single-corpus alignment dependence.}
\section{Threats to validity}
\TODO{leakage (addressed), contamination (addressed), fit collapse (guarded),
median vs mean reporting.}

% =====================================================================
\chapter{Future Work}\label{ch:future}
\section{Phase 2: intonation and vibrato}
\TODO{f0 extraction from monophonic audio; new channels; datasets (URMP, singing).}
\section{Phase 3: a differentiable-synthesizer audio likelihood}
The intended step to audio: $x=\Phi(z)\mathbf a+\varepsilon$, amplitudes
marginalised analytically, nonlinear positions inferred approximately.
\TODO{status: not yet started.}
\section{Geometric direction: the score-bundle connection Laplacian}
The current graph prior is the \emph{flat} case of a graph connection Laplacian;
non-trivial edge transforms give a genuine fibre bundle with measurable
curvature/holonomy. \TODO{see docs/geometry\_direction.md; litmus test.}
\section{Unknown score support}
\TODO{toward full AMT via a discrete latent layer.}

% =====================================================================
\chapter{Conclusion}\label{ch:conclusion}
\TODO{restate the validated contribution (structure + calibration, each
ingredient isolated), the boundary, and the road to audio.}

% =====================================================================
\appendix
\chapter{Full formulation}\label{app:formulation}
\TODO{\texttt{\textbackslash input\{formulation\}} here, or merge
docs/formulation.tex into this appendix.}

\chapter{Additional results}\label{app:extra}
\TODO{per-channel tables, ablations, reliability diagrams / PIT histograms.}

% =====================================================================
\begin{thebibliography}{99}
\bibitem{lindgren2011} F.~Lindgren, H.~Rue, J.~Lindstr\"om.
  An explicit link between Gaussian fields and Gaussian Markov random fields:
  the SPDE approach. \emph{JRSS-B}, 2011.
\bibitem{rueheld2005} H.~Rue, L.~Held.
  \emph{Gaussian Markov Random Fields: Theory and Applications}. 2005.
\bibitem{borovitskiy2021} V.~Borovitskiy et al.
  Mat\'ern Gaussian Processes on Graphs. \emph{AISTATS}, 2021.
\bibitem{ddsp} J.~Engel, L.~Hantrakul, C.~Gu, A.~Roberts.
  DDSP: Differentiable Digital Signal Processing. \emph{ICLR}, 2020.
\bibitem{musictransformer} C.-Z.~A.~Huang et al.
  Music Transformer. \emph{ICLR}, 2019.
\bibitem{remi} Y.-S.~Huang, Y.-H.~Yang.
  Pop Music Transformer (REMI). \emph{ACM MM}, 2020.
\bibitem{oore2018} S.~Oore et al.
  This Time with Feeling: Learning Expressive Musical Performance. 2018.
\bibitem{nanogpt} A.~Karpathy. nanoGPT. \url{https://github.com/karpathy/nanoGPT}.
\bibitem{cancino2018} C.~Cancino-Chac\'on et al.
  Computational Models of Expressive Music Performance: A Review.
  \emph{Front.\ Digit.\ Humanit.}, 2018.
\bibitem{asap} F.~Foscarin et al.
  ASAP: a dataset of aligned scores and performances. \emph{ISMIR}, 2020.
\bibitem{maestro} C.~Hawthorne et al.
  Enabling Factorized Piano Music Modeling with the MAESTRO Dataset.
  \emph{ICLR}, 2019.
\bibitem{vienna4x22} W.~Goebl. The Vienna 4x22 Piano Corpus. \TODO{cite details.}
\bibitem{TODO} \TODO{add remaining references (Singer--Wu connection Laplacian,
  Tymoczko, ATEPP, Aria-MIDI, alignment methods, \dots).}
\end{thebibliography}

\end{document}
